\input{TD.sty}
\newcommand{\pcsi}[1]{
}
\newcommand{\etudiant}[2][*]{
  \begin{itemize}
\ifthenelse{\equal{\detokenize{#1}}{\detokenize{*}}}
{\item \color{red}\fbox{\uppercase{\textbf{#2}}}~:}
{\item \fbox{\uppercase{\textbf{#2}}}~:}
\end{itemize}
}

\newcommand{\pc}[1]{
}

\begin{document}
\begin{center}
  \Large
  \textsc{TIPE de physique}
\end{center}

\etudiant{BIZID}
Effet Magnus : sujet riche, possibilité d'expérience et de numérique. On étudie l'effet en PC ; en tout cas la théorie est abordable. Par contre, lien avec le thème de l'année cycle, boucles ?

\etudiant{BRESCON}
Trompette : acoustique étudiée en PC, beaucoup d'expériences possibles. Sujet solide. Mais lien avec le thème de l'année cycle, boucles ?

\etudiant{BRUNO}
Horloge à balancier : sujet dans le thème de l'année. On semble tourner en rond, peut-être aller plus loin dans l'étude du mécanisme, raffiner le modèle (aller au-delà du pendule simple), élargir à d'autres types d'horloges. Mais sujet viable.

\etudiant{BUETAS} 
Cycles thermodynamiques. Sujet dans le thème de l'année. Peu de résultats pour l'instant apparemment, mais on peut élargir à tout type de cycle (moteur, réfrigérateur, autre). Sujet viable, mais il va falloir bosser davantage...

\etudiant{DROIN}
Vortex. Vaguement dans le thème de l'année. Sujet casse-gueule (Duval a bien galéré dessus). Les techniques d'imagerie d'écoulements, c'est trop compliqué. Il vaudrait mieux hanger de sujet.

\etudiant{FAURE}
Hélice d'avion. Aérodynamique étudiée en PC (même si ce sont vraiment des bases). Mais on ne dispose pas d'appareil de précision au lycée (Pitot et soufflerie bof...). Je conseillerais de changer.  

\etudiant{GABRIEL}
Bulles de champagne. Problème très intéressant et bien documenté. On peut faire des expériences simples et amusantes. La méca flu de PC permet déjà de dire beaucoup de choses. Lien avec le thème de l'année un peu tiré par les cheveux.... On peut garder le sujet.

\etudiant{GUELA} 
Boxe. Lien avec cycle, boucles ?? Peu de résultats, je sens qu'on va tourner en rond. Je conseille de changer de sujet.

\etudiant{LARROQUE - SAINT-PAUL} 
Mascara. Ok pour le thème de l'année (cils bouclés j'imagine ?) Travail déjà solide. Tension superficielle pas au programme malheureusement, mais littérature abondante, et beaucoup d'expériences possibles. Peut garder le sujet.

\etudiant{RUELLAND} 
Cycle réfrigérant. Sujet OK. Ne pas limiter au stirling, et changer de problématique, sinon on risque de tourner en rond. Ok pour garder le sujet, mais il faudra y réfléchir davantage.

\etudiant{SCHMITT}
Résistance au cycle de fonte. Sujet dans le thème de l'année. Résultats solides, c'est bien. Il faudra retravailler un peu la présentation des résultats, mais déjà beaucoup de travail. On peut garder le sujet.

\end{document}



